\section*{Capítulo \ref{ch:g12:matrix} --- Matrizes e grafos}

\begin{solution}{exo:g12:matrix:1}
\[
A + B = \begin{pmatrix} 3 & 2\\ 1 & 2\end{pmatrix}, \quad
AB = \begin{pmatrix} 4 & 2\\ 1 & 1\end{pmatrix}, \quad
BA = \begin{pmatrix} 2 & 4\\ 1 & 3\end{pmatrix}, \quad
A^2 = \begin{pmatrix} 1 & 4\\ 0 & 1\end{pmatrix}.
\]
Observe que $AB \neq BA$.
\end{solution}

\begin{solution}{exo:g12:matrix:2}
$\det A = 6 - 5 = 1 \neq 0$:
$A^{-1} = \begin{pmatrix} 3 & -5\\ -1 & 2 \end{pmatrix}$.
$\det B = 12 - 12 = 0$: $B$ não é \omterm{def:g12:matrix:inverse}{invertível}.
\end{solution}

\begin{solution}{exo:g12:matrix:3}
O sistema é $AX = Y$ com $A$ como no \cref{exo:g12:matrix:2} e
$Y = \begin{pmatrix} 1\\ 2\end{pmatrix}$:
\[
X = A^{-1}Y = \begin{pmatrix} 3 & -5\\ -1 & 2\end{pmatrix}
\begin{pmatrix} 1\\ 2\end{pmatrix}
= \begin{pmatrix} -7\\ 3\end{pmatrix}:
\qquad x = -7,\ y = 3 .
\]
\end{solution}

\begin{solution}{exo:g12:matrix:4}
\emph{1.} $N^2 = \begin{pmatrix} 0&1\\0&0\end{pmatrix}
\begin{pmatrix} 0&1\\0&0\end{pmatrix} = 0$, e $I$ comuta com toda
\omterm{def:g12:matrix:matrix}{matriz}.

\emph{2.} Como os dois termos comutam, o teorema binomial se aplica e
todos os termos que contêm $N^2$ se anulam:
\[
A^k = (2I + N)^k = 2^k I + k\,2^{k-1} N
= \begin{pmatrix} 2^k & k\,2^{k-1}\\ 0 & 2^k \end{pmatrix}.
\]
Verificação para $k = 2$:
$A^2 = \begin{pmatrix} 2&1\\0&2\end{pmatrix}^2
= \begin{pmatrix} 4&4\\0&4\end{pmatrix}$, e a fórmula dá
$2^2 = 4$, $2 \times 2 = 4$. \checkmark
\end{solution}

\begin{solution}{exo:g12:matrix:5}
\emph{Indução.} Para $n = 1$:
$A^1 = \begin{pmatrix} 0&1\\1&1\end{pmatrix}
= \begin{pmatrix} F_0 & F_1\\ F_1 & F_2\end{pmatrix}$. Suponha a fórmula
válida para $n$; então
\[
A^{n+1} = A^n A
= \begin{pmatrix} F_{n-1} & F_n\\ F_n & F_{n+1}\end{pmatrix}
\begin{pmatrix} 0 & 1\\ 1 & 1\end{pmatrix}
= \begin{pmatrix} F_n & F_{n-1} + F_n\\ F_{n+1} & F_n + F_{n+1}\end{pmatrix}
= \begin{pmatrix} F_n & F_{n+1}\\ F_{n+1} & F_{n+2}\end{pmatrix}.
\]
\emph{Identidade.} Para \omterm{def:g12:matrix:matrix}{matrizes} $2\times2$, o desenvolvimento mostra que
$\det(MN) = \det M \det N$; logo
$\det(A^n) = (\det A)^n = (-1)^n$, e
$\det A^n = F_{n-1}F_{n+1} - F_n^2$. (Esta é a \emph{identidade de
Cassini}.)
\end{solution}

\begin{solution}{exo:g12:matrix:6}
\emph{1.} Ordenando os vértices $1, 2, 3$:
\[
M = \begin{pmatrix} 0&1&1\\ 0&0&1\\ 1&0&0\end{pmatrix}, \quad
M^2 = \begin{pmatrix} 1&0&1\\ 1&0&0\\ 0&1&1\end{pmatrix}, \quad
M^3 = \begin{pmatrix} 1&1&1\\ 1&0&1\\ 1&0&1 \end{pmatrix}.
\]

\emph{2.} $\bigl(M^3\bigr)_{11} = 1$: exatamente um passeio fechado de
comprimento $3$ no vértice $1$, a saber $1 \to 2 \to 3 \to 1$. (O passeio
$1 \to 3 \to 1$ tem comprimento apenas $2$, e $1 \to 3$, depois
$3\to1$, depois $1\to3$ termina em $3$.)
\end{solution}

\begin{solution}{exo:g12:matrix:7}
\emph{1.} $a_{n+1} = 0.8a_n + 0.3b_n$, $b_{n+1} = 0.2a_n + 0.7b_n$:
$M = \begin{pmatrix} 0.8 & 0.3\\ 0.2 & 0.7\end{pmatrix}$.

\emph{2.} $MX = X$ dá $0.8a + 0.3b = a$, \emph{isto é}, $0.3b = 0.2a$,
logo $b = \frac23 a$; com $a + b = 1$: $a = 0.6$, $b = 0.4$.

\emph{3.} Usando $b_n = 1 - a_n$:
$a_{n+1} = 0.8a_n + 0.3(1 - a_n) = 0.5a_n + 0.3$, de modo que
\[
c_{n+1} = a_{n+1} - 0.6 = 0.5a_n + 0.3 - 0.6 = 0.5(a_n - 0.6) = 0.5\,c_n .
\]
Assim, $c_n = 0.5^n c_0 \to 0$: $a_n \to 0.6$ e $b_n \to 0.4$, qualquer
que seja a \omterm{def:g11:prob:rv}{distribuição} inicial.
\end{solution}

\begin{solution}{exo:g12:matrix:8}
\emph{1.} $\det P = -2$, logo
$P^{-1} = -\frac12\begin{pmatrix} -1 & -1\\ -1 & 1\end{pmatrix}
= \frac12\begin{pmatrix} 1 & 1\\ 1 & -1\end{pmatrix}$. Então
\[
AP = \begin{pmatrix} 4 & 2\\ 4 & -2 \end{pmatrix}, \qquad
D = P^{-1}AP = \frac12\begin{pmatrix} 1&1\\1&-1\end{pmatrix}
\begin{pmatrix} 4&2\\4&-2\end{pmatrix}
= \begin{pmatrix} 4 & 0\\ 0 & 2\end{pmatrix}.
\]

\emph{2.} De $A = PDP^{-1}$, uma indução imediata dá
$A^n = PD^nP^{-1}$ com
$D^n = \begin{pmatrix} 4^n & 0\\ 0 & 2^n\end{pmatrix}$, logo
\[
A^n = P D^n P^{-1}
= \begin{pmatrix} 4^n & 2^n\\ 4^n & -2^n\end{pmatrix}\cdot
\frac12\begin{pmatrix} 1&1\\1&-1\end{pmatrix}
= \frac12\begin{pmatrix} 4^n + 2^n & 4^n - 2^n\\
4^n - 2^n & 4^n + 2^n\end{pmatrix}.
\]
Aplicar $A^n$ a $X_0 = \begin{pmatrix} u_0\\v_0\end{pmatrix}$ reproduz
exatamente as fórmulas do \cref{ex:g12:matrix:coupled}.
\end{solution}

\begin{solution}{pb:g12:matrix:1}
\textbf{1.} $AB = \begin{pmatrix} 2 & 1\\ 4 & 3\end{pmatrix}$
e $BA = \begin{pmatrix} 3 & 4\\ 1 & 2\end{pmatrix}$: a
multiplicação de \omterm{def:g12:matrix:matrix}{matrizes} não é comutativa --- $B$ troca as
colunas quando está à direita e as linhas quando está à
esquerda.

\textbf{2.} \omterm{def:g11:vect:det}{Determinante} $6 - 5 = 1$: inversa
$\begin{pmatrix} 3 & -1\\ -5 & 2\end{pmatrix}$. Aplicando-a a
$\binom{4}{7}$: $x = 12 - 7 = 5$, $y = -20 + 14 = -6$.

\textbf{3.} $N^2 = 0$. Então $(I + N)^n = I + nN$ por indução:
$(I + nN)(I + N) = I + (n+1)N + nN^2 = I + (n+1)N$.

\textbf{4.} $A = \begin{pmatrix} 0&1&1\\ 1&0&1\\ 1&1&0
\end{pmatrix}$, e $A^3$ tem entradas diagonais iguais a $2$: de
cada vértice partem exatamente dois passeios fechados de
comprimento $3$ (o triângulo percorrido em um sentido ou no
outro) --- o teorema de contagem em ação.

\textbf{5.} $D^n = \begin{pmatrix} 2^n & 0\\ 0 & 2^{-n}
\end{pmatrix}$: uma direção explode, a outra morre --- destinos
diagonais são \omterm{def:g12:seq:sequence}{sequências} geométricas independentes.

\textbf{6.} $F^2 = \begin{pmatrix} 2 & 1\\ 1 & 1\end{pmatrix}$,
$F^3 = \begin{pmatrix} 3 & 2\\ 2 & 1\end{pmatrix}$,
$F^4 = \begin{pmatrix} 5 & 3\\ 3 & 2\end{pmatrix}$: Fibonacci
por toda parte; a conjectura é a do enunciado.

\textbf{7.} Se $F^n = \begin{pmatrix} F_{n+1} & F_n\\ F_n &
F_{n-1}\end{pmatrix}$, então
\[
F^{n+1} = F^n F =
\begin{pmatrix} F_{n+1} + F_n & F_{n+1}\\
F_n + F_{n-1} & F_n \end{pmatrix}
= \begin{pmatrix} F_{n+2} & F_{n+1}\\ F_{n+1} & F_n
\end{pmatrix} :
\]
a hereditariedade; o caso inicial $n = 1$ é o próprio $F$, com
a convenção $F_0 = 0$ (que estende a recorrência para trás).

\textbf{8.} $\det F = -1$, logo
$\det(F^n) = (\det F)^n = (-1)^n$; e, diretamente,
$\det(F^n) = F_{n+1}F_{n-1} - F_n^2$: Cassini, em uma linha. (A
regra do produto para \omterm{def:g11:vect:det}{determinantes} $2 \times 2$ é um
desenvolvimento agradável de cinco minutos.)

\textbf{9.} Canto superior direito de $F^m F^n$:
$F_{m+1}F_n + F_m F_{n-1}$; canto superior direito de
$F^{m+n}$: $F_{m+n}$. Para $m = n = 3$:
$F_4 F_3 + F_3 F_2 = 3 \times 2 + 2 \times 1 = 8 = F_6$.

\textbf{10.} Para $k = 1$: trivial. Se $F_n \mid F_{kn}$, a
fórmula de adição com $m = kn$ dá
$F_{(k+1)n} = F_{kn+1}F_n + F_{kn}F_{n-1}$: os dois termos são
múltiplos de $F_n$. Logo $F_n \mid F_{kn}$ para todo $k$:
confira que $F_3 = 2$ \omterm{def:g12:arith:divides}{divide} $F_6 = 8$ e $F_9 = 34$.

\textbf{11.} $F^{100} = F^{64} F^{32} F^4$: sete elevações ao
quadrado ($F^2, F^4, \dots, F^{64}$) mais duas combinações ---
nove multiplicações em vez de noventa e nove. É o truque da
tabela de dobros dos escribas egípcios, transposto dos números
para as \omterm{def:g12:matrix:matrix}{matrizes}: escreva $100$ em binário e multiplique os
dobros de que precisa.

\textbf{12.} $0.8 + 0.2 = 1$ e $0.4 + 0.6 = 1$: amanhã terá
\emph{algum} tempo --- cada coluna é uma \omterm{def:g11:prob:rv}{distribuição} de
\omterm{def:g10:proba:distribution}{probabilidade} completa, de modo que as \omterm{def:g10:proba:distribution}{probabilidades} se
conservam.

\textbf{13.} Amanhã: $\binom{0.8}{0.2}$. Depois de amanhã:
$M\binom{0.8}{0.2} = \binom{0.72}{0.28}$.

\textbf{14.} $Mv = v$ com $v = \binom{s}{c}$, $s + c = 1$:
$0.8s + 0.4c = s$ dá $0.4c = 0.2s$, $s = 2c$:
$v = \binom{2/3}{1/3}$. A longo prazo, dois dias em cada três
são de sol --- não importa como esteja hoje.

\textbf{15.} A partir de $\binom01$: componentes de sol $0.4$,
$0.56$, $0.624$, $0.6496$; distâncias a $\frac23$: $0.267$,
$0.107$, $0.043$, $0.017$ --- cada passo multiplica a
diferença por exatamente $0.4$ (o segundo autovalor da
máquina): convergência geométrica para o regime estacionário.

\textbf{16.} Colunas (a partir de $A$, $B$, $C$):
$P = \begin{pmatrix} 0 & 0 & 1\\ \frac12 & 0 & 0\\
\frac12 & 1 & 0\end{pmatrix}$. Regime estacionário:
$v_A = v_C$, $v_B = \frac{v_A}{2}$,
$v_C = \frac{v_A}{2} + v_B$; somando $1$:
$v = \left(\frac25, \frac15, \frac25\right)$. Ordenação: $A$ e
$C$ empatam em primeiro, $B$ fica em último.

\textbf{17.} $C$ recebe \emph{todo} o tráfego de $B$ e metade
do de $A$, e devolve tudo a $A$: o regime estacionário mede
onde o navegante \emph{passa o tempo}, não quantos links
apontam para a página --- um link vindo de uma página popular
pesa mais que vários vindos de páginas desertas. Essa
ponderação recursiva é exatamente a ideia fundadora do Google;
o amortecimento cuida das armadilhas e dos becos sem saída.

\textbf{18.} $A^n = P D^n P^{-1}$ dá
$u_n = \frac{4^n + 2^n}{2}$ (e
$v_n = \frac{4^n - 2^n}{2}$). Verificação: $u_1 = 3$,
$u_2 = 10$, $u_3 = 36$; diretamente:
$(1,0) \to (3,1) \to (10,6) \to (36, 28)$: bate.

\textbf{19.} A diagonalização passa para \omterm{def:g10:coordgeom:system}{coordenadas} nas quais
o sistema acoplado se desfaz em \omterm{def:g12:seq:sequence}{sequências} geométricas
independentes --- cada autovalor corre sua própria corrida. O
regime estacionário de Markov é o autovetor associado ao
autovalor $1$, e a taxa de convergência da questão 15 é o
autovalor seguinte: a máquina do tempo era, desde o começo,
uma história de autovetores.

\textbf{20.} Contabilidade: um sistema é uma única \omterm{def:g10:algebra:equation}{equação}
matricial, resolvida por uma única inversa. Combinatória: as
potências da \omterm{def:g12:matrix:graph}{matriz de adjacência} contam passeios, links,
conexões. Evolução: as potências da máquina levam os estados
a seus destinos, e as direções próprias (a direção áurea de
Fibonacci, o regime estacionário do tempo, o \omterm{def:g11:vect:vector}{vetor} de
ordenação da web) são esses destinos. A álgebra linear, nos
volumes de graduação, é a ciência exatamente disso.
\end{solution}
